\lecture{2}
\title{ERM and MLE}

\begin{document}

Let's still assume all data is drawn (i.i.d.) from some fixed distribution $p$,
but say we don't know what $p$ is.

\begin{remark}
    This isn't entirely realistic---sensor data, say, might have temporal correlations.
\end{remark}

\subsection{Empirical Risk Minimization}

Here's a very simple idea:
Use our training data $\{(x^{(n)}, y^{(n)})\}_{n = 1}^N$ to
derive an \textbf{empirical distribution} like so:
\[ \widehat{p}(x, y) := \frac{1}{N} \sum_{n = 1}^N \delta_{(x^{(n)}, y^{(n)})}(x, y). \]
Here, $\delta_k(x)$ is the Dirac delta function, defined so that
$\int_{\mathbb{R}^D} \delta_k(x) \, \mathrm{d}x = 1$ and $\delta_k(x) = 0$ for all $x \neq k$.

Using the empirical distribution also gives us an \textbf{empirical risk} that looks like:
\[ \widehat{\text{risk}}(h) = \frac{1}{N} \sum_{n = 1}^N L(y^{(n)}, h(x^{(n)})). \]
And so we might choose $h$ according to \textbf{Empirical Risk Minimization (ERM)}, like so:
\[ h_{\textsc{Erm}} = \argmin_{h \in \mathcal{H}} \widehat{\text{risk}}(h), \]
where $\mathcal{H}$ is some class of candidate decision rules.

\subsection{Overfitting}

Of course, this stupid solution fails in stupid ways.

\begin{quote}
    \textbf{Example.} Given a training set of spam / not-spam emails,
    decide whether future emails are spam.

    \textbf{ERM Solution.} An email is declared spam if its timestamp
    matches that of a spam email in the training set.
\end{quote}

This is an example of \textbf{overfitting}.  Here are two approaches by which we might fix it:
\begin{enumerate}
    \item Restrict the hypothesis class $\mathcal{H}$.
    \item Approximate $p(x, y)$ differently.
\end{enumerate}
Let's take the latter of the two approaches for now: how can we better estimate $p(x, y)$?

\subsection{Maximum Likelihood Estimation}

\begin{remark}
    This section is written with the assumption of 18.650,
    particularly up through \href{/18.650/lectures/9}{Lecture 9},
    which covers the MLE.
\end{remark}

Rather than naively approximating $p(x, y)$ using the empirical distribution $\widehat{p}(x, y)$,
let's suppose $p$ belongs to a \emph{parametric model}
and estimate the parameter of $p$ via the \emph{MLE}.

\begin{problem}{Bernoulli MLE}
    Suppose, for simplicity, that we only care about
    approximating the distribution of labels $\{y^{(n)}\}_{n = 1}^N$
    and ignore the features $\{x^{(n)}\}_{n = 1}^N$.
    Also assume that the $y^{(n)} \in \{0, 1\}$
    are sampled from $\mathrm{Ber}(\theta)$ for some $\theta \in [0, 1]$.

    As a function of the samples $\{y^{(n)}\}_{n = 1}^N$,
    determine the MLE $\hat{\theta}$.
\end{problem}

\begin{solution}
    Because the labels $\{y^{(n)}\}_{n = 1}^N$ are independent, the likelihood is:
    \[ p(\mathcal{D} \mid \theta)
    = \prod_{n = 1}^N p(y^{(n)} \mid \theta)
    = \prod_{n = 1}^N \theta^{y^{(n)}} (1 - \theta)^{1 - y^{(n)}}. \]
    We'd like to maximize the above,
    so it is natural to differentiate the above with respect to $\theta$\dots
    but the result seems very messy.
    The trick, of course, is to consider the log-likelihood:
    \[ \ell(\theta) = \log p(\mathcal{D} \mid \theta)
    = \sum_{n = 1}^N \left[
        y^{(n)} \log \theta + (1 - y^{(n)}) \log (1 - \theta)
    \right]. \]
    Technically we should treat $\theta \in \{0, 1\}$ separately.
    That aside, we can now differentiate $\ell(\theta)$ easily.
    \[ \frac{d}{d\theta} \ell(\theta)
    = \frac{1}{\theta}\left[\sum_{n = 1}^N y^{(n)}\right]
    - \frac{1}{1 - \theta}\left[\sum_{n = 1}^N (1 - y^{(n)})\right]. \]
    The above equals $0$ exactly at $\hat{\theta} = \frac{1}{N}\sum_{n = 1}^N y^{(n)}$.
    It now remains to check that this critical point is a maximum:
    \[ \frac{d^2}{d\theta^2} \ell(\theta)
    = -\frac{1}{\theta^2}\left[\sum_{n = 1}^N y^{(n)}\right]
    - \frac{1}{(1 - \theta)^2}\left[\sum_{n = 1}^N (1 - y^{(n)})\right] < 0. \]
    So the MLE must be $\hat{\theta} = \frac{1}{N}\sum_{n = 1}^N y^{(n)}$. ~ $\blacksquare$
\end{solution}

\subsection{MLE Advantages and Disadvantages}

The MLE has some advantages!
\begin{itemize}
    \item \textbf{Computation.}
    As long as the likelihood is differentiable,
    it is often very easy and fast to computationally find the MLE.
    \item \textbf{Reparametrization.}
    The MLE is invariant to reparametrization
    by a bijective function $\eta = f(\theta)$.
    (So nothing unexpected changes if we use, say,
    $\sigma^2$ as a parameter instead of $\sigma$.)
\end{itemize}

But it also has some disadvantages.
\begin{itemize}
    \item \textbf{No Indication of Certainty.}
    An MLE computed using ten datapoints
    looks the same as an MLE computed using ten million.
    \item \textbf{Poor Performance for Small $N$.}
    If we sample from $\mathrm{Ber}(0.2)$,
    it's likely the MLE will be stuck at $\hat{\theta} = 0$ for a while\dots
    \item \textbf{Likelihood Inflation.}
    Consider a \emph{mixture model} of two Gaussians
    $\mathcal{N}(\mu_0, \sigma_0^2)$ and $\mathcal{N}(\mu_1, \sigma_1^2)$
    with mixture weight $\lambda$.
    The MLE has the freedom to decide the values of $(\mu_0, \sigma_0^2, \mu_1, \sigma_1^2)$
    however it wants\dots so what if it does this?
    \begin{quote}
        \textbf{MLE Proposal.} Estimate $\hat{\mu}_0 = y^{(1)}$ and $\hat{\sigma}_0^2 \ll 1$.
        Don't even bother with $\hat{\mu}_1$ and $\hat{\sigma}_1^2$.
    \end{quote}
    With this proposal, the likelihood looks like this:
    \[ p(\mathcal{D} \mid \mu_0, \sigma_0^2, \mu_1, \sigma_1^2)
    \geq \left[(1 - \lambda) \cdot \frac{1}{\sqrt{2\pi \sigma_0^2}}\right]
    \cdot \prod_{n = 2}^N \left[
        \lambda \cdot \frac{1}{\sqrt{2\pi \sigma_1^2}} \cdot
        \exp\left(-\frac{(y^{(n)} - \mu_1)^2}{2\sigma_1^2}\right)
    \right]. \]
    For choices of $\sigma_0^2$ sufficiently small,
    the MLE can have as big a likelihood as it wants!
    But even though it's maximizing the likelihood,
    picking an estimator by only looking at a single sample $y^{(1)}$
    does not seem very helpful\dots
\end{itemize}

\end{document}